\section{Discussion}
\label{sec:discussion}

\subsection{Consolidated results}

\begin{table}[htbp]
\centering
\caption{Headline results. Every value is read from the simulation output at
compile time, so this table cannot drift out of step with the code.}
\label{tab:summary}
\small
\begin{tabular}{llr}
\toprule
Phase & Quantity & Value \\
\midrule
Transfer   & Total impulsive $\dv$              & $\mDVtot\unit{m/s}$ \\
           & Phasing wait                       & $\mTwaitH\unit{h}$ \\
           & Mission duration                   & $\mTmissionH\unit{h}$ \\
           & Numerical vs analytical error      & $\mKeplerErrMm\unit{mm}$ \\
\midrule
Proximity  & Injection error                    & $\mDrZero\unit{m}$, $\mDvZero\unit{m/s}$ \\
           & Uncorrected drift, 3 orbits        & $\mDriftKm\unit{km}$ \\
           & Hold acquisition $\dv$             & $\mDVhold\unit{m/s}$ \\
           & Docking $\dv$                      & $\mDVdock\unit{m/s}$ \\
           & Linearisation error, 3 orbits      & $\mNlSep\unit{m}$ \\
\midrule
Perturbed  & $J_2$ miss, no retargeting         & $\mMissJtwo\unit{km}$ \\
           & Third-body contribution            & $\mMissTB\unit{m}$ \\
           & Mid-course correction budget       & $\mDVmcc\unit{m/s}$ \\
\midrule
CR3BP      & Converged miss                     & $\mCRmiss\unit{m}$ \\
           & Correction to the two-body design  & $\mCRextra\unit{mm/s}$ \\
\midrule
Robotics   & Reach, 5R                          & $\mReach\unit{m}$ \\
           & Peak joint torque at $100\unit{N}$ & $\mTauMax\unit{N\,m}$ \\
           & IK position residual               & $\mIKres\unit{mm}$ in $\mIKiters$ iterations \\
\bottomrule
\end{tabular}
\end{table}

\subsection{Verification summary, including what disagrees}
\label{sec:discussion:verif}

An automated audit re-derives every design scalar in closed form from the raw
constants and compares the rest against an independent implementation of the
same mission geometry. It records $60$ passes, one warning and no failures. The
three results worth discussing are the ones that did not match exactly.

\paragraph{Docking $\dv$: $\mDVdock$ against $0.1567\unit{m/s}$ ($0.02\%$).}
The forced \vbar{} profile is almost pure geometry, so agreement at this level
means the leg targeting, waypoint spacing and braking impulse are modelled
identically. This is the strongest single confirmation that the proximity module
is faithful.

\paragraph{Third body: $\mMissTB$ against $116.0\unit{m}$ ($0.2\%$).}
Agreement at this level on a $116\unit{m}$ effect tests two things that are easy
to get wrong: the retention of the indirect term in~\eqref{eq:thirdbody}, and
the use of the synodic rate for the Earth ephemeris. Dropping the indirect term
inflates the number by orders of magnitude, so a match this tight is real
evidence rather than a coincidence.

\paragraph{$J_2$ miss: $\mMissJtwo\unit{km}$ against $14.56\unit{km}$ ($+4\%$).}
The four plausible causes were checked in the source rather than assumed: the
departure impulse is applied along the perturbed velocity, the target is
propagated with the same perturbed dynamics as the chaser, the full
three-component $J_2$ formula is used, and the Keplerian control case through
the identical code path closes to $5.6\times10^{-3}\unit{m}$. None is at fault.
The residual is sensitivity, not error. The miss is a \emph{difference} of two
nearly equal secular phase slips accumulated over ten and a half hours, so it
inherits the sensitivity of both to the exact burn epoch and to the integrator
tolerance. Two correct implementations agreeing to $4\%$ on such a quantity is
the expected outcome, and the discrepancy was documented rather than removed.

\paragraph{CR3BP correction: $\mCRextra\unit{mm/s}$ against $1.5\unit{mm/s}$.}
The single warning. As explained in Section~\ref{sec:cr3bp}, the planar shooting
problem is rank-deficient by one, so the zero-miss set is a curve and different
optimiser starts land at different points on it. Both figures support the same
operational statement -- Earth tides move this budget by millimetres per second
-- and forcing agreement would be curve-fitting rather than physics.

\paragraph{Hold $\dv$: $\mDVhold$ against $1.4\unit{m/s}$.}
Not comparable: the injection error is a random draw and the two runs use
different realisations. Fed the reference draw, our targeting returns
$1.419\unit{m/s}$, reproducing the reference to $1.4\%$. That test is also what
settles the frame convention of Section~\ref{sec:proximity:frame}.

\subsection{Limitations}
\label{sec:limits}

This is an engineering-faithful simulation of a restricted model. The following
would each change a specific number, and are named so that no reader has to
discover them.

\begin{itemize}
\item \textbf{Equatorial, coplanar orbits only.} No inclination difference and no
      relative right ascension, so the out-of-plane budget that dominates real
      rendezvous planning never appears. This is the single largest omission.
\item \textbf{Gravity stops at $J_2$.} The real lunar field is dominated by mass
      concentrations rather than a clean oblateness term \cite{zuber_grail}, and
      the $C_{22}$ ellipticity is comparable to $J_2$ for some geometries. Low
      lunar orbits are genuinely unstable over weeks under the true field.
      Nothing here runs that long, but the frozen-orbit design a real mission
      would use is a numerical search this model cannot support.
\item \textbf{No solar radiation pressure.} Estimated at $10^{-7}\unit{m/s^2}$
      and excluded with the reasoning given in Section~\ref{sec:perturbations};
      the flag exists but no cannonball model is wired in.
\item \textbf{Impulsive burns.} Finite-burn losses, throttle profiles, attitude
      slews and propellant slosh are all absent. For a $60\unit{m/s}$ impulse the
      gravity loss is small, but it is not zero.
\item \textbf{Circular target orbit, not a halo.} A realistic Artemis-era
      architecture places the mothership in a near-rectilinear halo orbit, where
      the rendezvous geometry and the phasing analysis are qualitatively
      different problems.
\item \textbf{The arm is kinematic and static.} Five degrees of freedom, a base
      held fixed by assumption, no joint flexibility, no contact dynamics, and no
      free-floating momentum coupling. The local frame is treated as inertial
      during berthing.
\item \textbf{One random draw.} The injection error is a single realisation, not
      a Monte Carlo campaign. The proximity $\dv$ figures are therefore one
      sample of a distribution whose spread is not characterised.
\end{itemize}

\subsection{Conclusions}

Three results are worth carrying away.

First, the propellant cost of this return phase is dominated by the transfer,
not by the corrections: $\mDVtot\unit{m/s}$ for the Hohmann climb against
$\mDVprox\unit{m/s}$ for all proximity operations, roughly $1\%$. What the
ascent dispersion actually costs is small, provided it is corrected early.

Second, the binding constraint is time, not propellant. The transfer lasts
$66$ minutes; waiting for the phasing window takes $\mTwaitH$ hours. Any
architecture that needs a fast return has to buy it with a non-Hohmann transfer,
and the CR3BP analysis shows that the extra energy will not come from
three-body effects at this altitude.

Third, the perturbation that matters at $100$--$400\unit{km}$ is lunar
oblateness, not the Earth. $J_2$ displaces the meeting point by
$\mMissJtwo\unit{km}$ and demands a $\mDVmcc\unit{m/s}$ mid-course correction;
the Earth contributes $\mMissTB\unit{m}$ and the full three-body treatment
changes the budget by $\mCRextra\unit{mm/s}$. Effort spent modelling the Earth
here is effort not spent on the lunar field, which is the opposite of the
intuition one brings from Earth orbit.

\subsection{Reproducibility}

The complete simulation runs unattended from a single command and takes about
ten minutes, of which roughly one is physics and the rest is rendering. It
requires base MATLAB only: no toolbox, no external ephemeris, no downloaded
imagery. All results in this report, including every figure, were produced by
the run whose audit table is reproduced in the repository as
\texttt{results/AUDIT\_REPORT.md}.
